\subsection{Benchmark Infrastructure}
\label{subsec::benchmark_infra}

The Cosmos benchmark system manages evaluation jobs for Cosmos models and stores both generated artifacts and evaluation results. An orchestration layer schedules generation, scoring, and endpoint evaluation jobs on Lepton or Slurm clusters and tracks the execution status of each stage. For every run, the system records metadata including the model checkpoint, code version, selected benchmarks, generation settings, benchmark-specific parameters, and associated datasets. Together, these records establish full traceability between each reported score and the exact model weights, inputs, parameter configurations, and evaluation code used to produce it.

The system supports the heterogeneous benchmark suite described in \cref{sec::results} without requiring all benchmarks to share a common implementation. Benchmarks evaluate generated video, audio, action trajectories, and text responses from reasoner models across a diverse set of criteria, including visual fidelity, audio quality, audio-visual synchronization, prompt and control adherence, action or trajectory accuracy, task completion, physical plausibility, and reasoning correctness. Evaluators include integrations with open-source libraries and public benchmark suites, as well as custom evaluators developed specifically for Cosmos. Scoring methods span reference-based error metrics, perceptual and temporal consistency measures, audio-video alignment metrics, VLM-based judges, human annotations, and exact-match or numeric-answer evaluation.

For Generator evaluation, benchmarking is separated into generation and scoring stages. Generation jobs execute models using either the PyTorch inference pipeline or one of the serving frameworks described in \cref{subsec::serving_infrastructure}, and write generated outputs to object storage. Scoring jobs subsequently consume these stored artifacts, compute per-sample and aggregate metrics, and record results together with run metadata. This decoupled design allows outputs to be rescored with new metrics or evaluators without rerunning generation.

For Reasoner evaluation, we use the VLMEvalKit~\citep{duan2024vlmevalkit} framework together with vLLM. These jobs send prompts and multimodal inputs to deployed model endpoints, process model responses, and record benchmark scores and associated metadata.

Evaluation scores and run metadata are stored in a relational database, while generated artifacts are stored in object storage. Human-evaluation annotations are stored together with the evaluated artifacts and question sets, and aggregate human-evaluation results are tracked alongside automated metrics. A benchmark portal provides access to these records through dashboards, leaderboards, example-level inspection tools, and model-to-model or checkpoint-to-checkpoint comparisons.
